% --------------------------------------------------------------------
%  FIGURE FRAMES
% --------------------------------------------------------------------
%  Three ways to put a drawing, chart or screenshot on the page:
%
%      \begin{utpfigure} ... \end{utpfigure}
%          a centred frame with no caption, for a diagram the text
%          talks about directly. An optional title puts a bar on top:
%          \begin{utpfigure}[Where the VLAN tag is added] ...
%
%      \begin{utptable} ... \end{utptable}
%          the same frame in teal, around a tabular or tabularx.
%
%      \begin{utpfig}[short entry]{The caption.}{fig:label}
%        \begin{tikzpicture} ... \end{tikzpicture}
%      \end{utpfig}
%          a frame with a numbered caption that \figref can point to.
%          The optional first argument is the short form for a list of
%          figures. A page never breaks between the drawing and its
%          caption.
%
%  A screenshot or slide lifted from the lecturer's material goes in a
%  titled frame, so the reader can tell it from a drawing made for the
%  notes; the optional width defaults to 0.9 and 0.78 of the line:
%
%      \sourceimage[0.8\linewidth]{Static-routing topology}{static-topology.png}
%      \compactsourceimage{The same, kept to a smaller height}{static-topology.png}
%
%  \fignote{...} adds a short muted line under a drawing, inside its frame.

\newtcolorbox{utpfigure}[1][]{utpplain,width=\linewidth,colback=white,colframe=UTPBlue!28, boxrule=0.7pt,arc=2mm, left=2.2mm,right=2.2mm,top=2.2mm,bottom=1.8mm, before skip=0.75em,after skip=0.85em, before upper={\centering}, halign=center, fonttitle=\bfseries\color{white},colbacktitle=UTPBlue,title={#1}}
\newtcolorbox{utptable}{utpplain,width=\linewidth,colback=white,colframe=UTPTeal!32, boxrule=0.7pt,arc=2mm, left=2.2mm,right=2.2mm,top=1.8mm,bottom=1.8mm, before skip=0.75em,after skip=0.85em, halign=center}
\newcommand{\utpbreakableframes}{%
  \renewtcolorbox{utpfigure}[1][]{utpplain,breakable,colback=white,colframe=UTPBlue!35, boxrule=0.5pt,arc=2mm,left=2.5mm,right=2.5mm,top=2.2mm,bottom=2.2mm, halign=center, fonttitle=\bfseries\color{white},colbacktitle=UTPBlue,title={##1}}%
  \renewtcolorbox{utptable}{utpplain,breakable,colback=white,colframe=UTPBlue!28, boxrule=0.45pt,arc=2mm,left=2.2mm,right=2.2mm,top=1.8mm,bottom=1.8mm}%
}

\captionsetup{
  hypcap=false,
  font={small},
  labelfont={bf,color=UTPNavy},
  textfont={color=UTPInk},
  labelsep=period,
  justification=centering,
  skip=4pt
}

\newenvironment{utpfig}[3][\relax]
  {\def\utpfigshort{#1}\def\utpfigcap{#2}\def\utpfiglab{#3}%
   \par\addvspace{0.55em}%
   \begin{tcolorbox}[utpplain,breakable,colback=white,colframe=UTPBlue!32,
     boxrule=0.5pt,arc=2mm,left=2.5mm,right=2.5mm,top=2.6mm,bottom=2.2mm]
   \centering}
  {\par\nopagebreak\vspace{0.45em}%
   % The caption is boxed: \captionof starts with a write and a zero
   % rule, and the glue after them would otherwise be a place where the
   % page could break, leaving the caption alone at the top of the next.
   % Inside the box the caption is told it sits below the drawing, which
   % it can no longer see, so the gap stays above it.
   \nointerlineskip
   \vbox{\setlength{\parskip}{0pt}\captionsetup{position=below}%
   \ifx\utpfigshort\relax
     \captionof{figure}{\utpfigcap}%
   \else
     \captionof{figure}[\utpfigshort]{\utpfigcap}%
   \fi
   \label{\utpfiglab}}%
   \end{tcolorbox}%
   \par\addvspace{0.55em}}

\newcommand{\sourceimage}[3][0.9\linewidth]{%
  \begin{utpfigure}[Source visual --- #2]%
    \includegraphics[width=#1,height=0.56\textheight,keepaspectratio]{#3}%
  \end{utpfigure}}
\newcommand{\compactsourceimage}[3][0.78\linewidth]{%
  \begin{utpfigure}[Source visual --- #2]%
    \includegraphics[width=#1,height=0.42\textheight,keepaspectratio]{#3}%
  \end{utpfigure}}

\newcommand{\fignote}[1]{\par\smallskip{\scriptsize\color{UTPMuted}#1\par}}
